% Generated from pinned audit evidence.
\newcommand{\TrainSeconds}{361.6}
\newcommand{\TrainMemory}{0.543}
\newcommand{\FirstIdentity}{1.17}
\newcommand{\FirstShuffled}{1.17}
\newcommand{\FirstZero}{0.00}
\newcommand{\BlockIdentity}{0.586}
\newcommand{\BlockGain}{0.098}
\newcommand{\BlockLow}{0.000}
\newcommand{\BlockHigh}{0.307}
\newcommand{\BlockFone}{0.320}
\newcommand{\RunDecision}{reject}

% Generated from the pinned V54 audit; no fitted presentation values.
\newcommand{\DensitySeconds}{207.8}
\newcommand{\DensityMemory}{0.400}
\newcommand{\DensityNll}{15.219}
\newcommand{\DensityShuffledNll}{15.254}
\newcommand{\DensityAdvantage}{0.0355}
\newcommand{\DensityLow}{-0.1014}
\newcommand{\DensityHigh}{0.1693}
\newcommand{\DensityIdentity}{1.562}
\newcommand{\DensityFone}{0.350}
\newcommand{\DensityCopyFone}{0.99949}
\newcommand{\DensityDecision}{reject}

% Generated from the locked V55 matched audit.
\newcommand{\JointNll}{7.026}
\newcommand{\ControlNll}{7.579}
\newcommand{\JointNllGain}{0.553}
\newcommand{\JointNllLow}{0.450}
\newcommand{\JointNllHigh}{0.665}
\newcommand{\JointTime}{342.4}
\newcommand{\ControlTime}{218.0}
\newcommand{\JointMemory}{0.648}
\newcommand{\ControlMemory}{0.400}
\newcommand{\JointSampleFone}{0.362}
\newcommand{\JointPrefixGain}{0.185}
\newcommand{\JointDecision}{reject}

% Generated from unchanged V56 report.
\newcommand{\BeamAccuracy}{11.621}
\newcommand{\BeamShuffled}{9.863}
\newcommand{\GreedyAccuracy}{8.252}
\newcommand{\BeamGain}{3.369}
\newcommand{\BeamGainLow}{2.351}
\newcommand{\BeamGainHigh}{4.412}
\newcommand{\BeamHistoryGain}{1.758}
\newcommand{\BeamHistoryLow}{0.966}
\newcommand{\BeamHistoryHigh}{2.552}
\newcommand{\BeamRepeat}{25.446}
\newcommand{\TargetRepeat}{0.112}
\newcommand{\BeamSeconds}{269.3}
\newcommand{\BeamMemory}{0.154}
\newcommand{\BeamDecision}{reject\_decoder}

% Generated from the completed V57 audit.
\newcommand{\QaExact}{3.906}
\newcommand{\QaFone}{15.718}
\newcommand{\QaBlankExact}{0.391}
\newcommand{\QaBlankFone}{5.553}
\newcommand{\QaControlExact}{0.000}
\newcommand{\QaControlFone}{6.582}
\newcommand{\QaQuestionGain}{10.165}
\newcommand{\QaQuestionLow}{7.043}
\newcommand{\QaQuestionHigh}{13.371}
\newcommand{\QaPairExact}{0.254}
\newcommand{\QaMotor}{100.000}
\newcommand{\QaSeconds}{278.572}
\newcommand{\QaMemory}{1.053}
\newcommand{\QaControlSeconds}{269.126}
\newcommand{\QaControlMemory}{1.049}
\newcommand{\QaDecision}{reject\_reading\_model}

\newcommand{\PaperTitle}{Learning conditional language\\in continuous image space}
\newcommand{\PaperSummary}{Written language includes historical, handwritten and regional forms that a fixed character inventory cannot faithfully represent. We investigate language modelling with writing images as input and output on one consumer GPU. A continuous image interface supports normalized raster likelihood and independent generation without a deployed character vocabulary. Joint reader--density learning improves target likelihood to \JointNll\ versus \ControlNll\ nats per image, while autonomous streams remain inaccurate. A law over complete observed image answers enables limited Chinese reading: 25 of 256 fixed development answers are exact, versus three with question pixels blanked and 21 after matched additional glyph-answer training. Raster F1 reaches 27.033\%, a 4.893-point gain [0.469, 9.047], while the exact-match difference remains unresolved. Both questions are correct in only four of 393 shared-paragraph pairs. Full adaptation of a 477M-parameter native pixel-language foundation improves held-out answer-pixel NLL from 0.6842 to 0.1548 nats but produces zero exact answers. Its independent runtime is reproducible, while free writing remains fragmented. Faithful forms and likelihood therefore require separate tests from useful language; a high-quality general image-based language model remains unresolved.}
